\documentclass[abstract=on,10pt,a4paper,bibliography=totocnumbered]{article}
\usepackage[paper=a4paper,left=35mm,right=35mm,top=25mm,bottom=30mm]{geometry}
\usepackage[doublespacing]{setspace}
\usepackage[english]{babel}
\usepackage[utf8]{inputenc}
\usepackage{amsmath}
\usepackage{colortbl}
\usepackage{amsfonts}
\usepackage{amssymb}
\usepackage{gensymb}
\usepackage{textcomp}
\usepackage{graphicx}
\usepackage{tikz}
\usepackage{enumerate}
\usepackage{enumitem}
\usepackage{subcaption}
\usepackage{booktabs}
\usepackage[hidelinks]{hyperref}
\usepackage[nameinlink]{cleveref}
\usepackage{multirow}
\usepackage{arydshln}
\usepackage[flushleft]{threeparttable}
\usepackage[nomarkers, nolists]{endfloat}
\usepackage{scalerel}
\usepackage{makecell}
\usepackage{ifthen}
\usepackage{tikz}
\usepackage{lineno}
\usepackage{adjustbox}
\usetikzlibrary{svg.path}

%------------------------------------------------------------------------------
%  Bibtex vs Biblatex
%------------------------------------------------------------------------------
% While bibtex is very nice to use, it's insanely slow. I therefore work with
% bibtex until the final stage. Set the value below to "true" or "false" to pick
% which one should be used.
\newboolean{usebiblatex}
\setboolean{usebiblatex}{true}

% Depending on the question, the commands below will prepare the document
% accordingly
\ifthenelse{\boolean{usebiblatex}}{
  %-----------------------------------------------------------------------------
  %  If using biblatex
  %-----------------------------------------------------------------------------
  % Load required packages and refer to correct bibliography
  \usepackage{csquotes}
  \usepackage[
      backend=biber
    , style=apa          % To display Author-year in the text
    , natbib=true        % Necessary to use citep, etc.
    , hyperref=true      % To have clickable links in the references
    , useeditor=false    % Don't add editor to publications
    , sorting=ynt        % Sort by year, name, and title
    , uniquename=false   % Avoid that initials are put
  ]{biblatex}
  \addbibresource{Literature.bib}

  % To include author in link
  \DeclareFieldFormat{citehyperref}{%
  \DeclareFieldAlias{bibhyperref}{noformat}% Avoid nested links
  \bibhyperref{#1}}
  \DeclareFieldFormat{textcitehyperref}{%
  \DeclareFieldAlias{bibhyperref}{noformat}% Avoid nested links
  \bibhyperref{%
    #1%
    \ifbool{cbx:parens}
      {\bibcloseparen\global\boolfalse{cbx:parens}}
      {}}}
  \savebibmacro{cite}
  \savebibmacro{textcite}
  \renewbibmacro*{cite}{%
    \printtext[citehyperref]{%
      \restorebibmacro{cite}%
      \usebibmacro{cite}}}
  \renewbibmacro*{textcite}{%
    \ifboolexpr{
      ( not test {\iffieldundef{prenote}} and
        test {\ifnumequal{\value{citecount}}{1}} )
      or
      ( not test {\iffieldundef{postnote}} and
        test {\ifnumequal{\value{citecount}}{\value{citetotal}}} )
    }
      {\DeclareFieldAlias{textcitehyperref}{noformat}}
      {}%
    \printtext[textcitehyperref]{%
      \restorebibmacro{textcite}%
      \usebibmacro{textcite}}}

  % Remove editor form everything but books
  \AtEveryBibitem{%
   \ifentrytype{book}{}{%
    \clearname{editor}
   }
  }
}{
  %-----------------------------------------------------------------------------
  %  If using bibtex
  %-----------------------------------------------------------------------------
  % Load required packages and select citation style
  \usepackage[round]{natbib}
  \bibliographystyle{apalike}
}

%------------------------------------------------------------------------------
%	Some Styling
%------------------------------------------------------------------------------
% Link colors
\definecolor{linkcolor}{HTML}{000000}
\hypersetup{
  colorlinks = true,
  linkcolor  = linkcolor,
  urlcolor   = linkcolor,
  citecolor  = linkcolor,
}

% Changing the style of captions in figures etc.
\captionsetup{labelfont=bf, format=plain, font=small}

% Change how equations are referenced
\renewcommand{\theequation}{Equation \arabic{equation}}%

% Avoid skip when including text from external .tex file
\newcommand{\inputy}[1]{\input{#1}\unskip}

% For supplementary material
\newcommand{\beginappendix}{%
  \setcounter{table}{0}
  \renewcommand{\thetable}{S\arabic{table}}%
  \setcounter{figure}{0}
  \renewcommand{\thefigure}{S\arabic{figure}}%
  \setcounter{equation}{0}
  \renewcommand{\theequation}{Equation S\arabic{equation}}%
  \setcounter{section}{0}
  \renewcommand{\thesection}{A.\arabic{section}}%
}

%------------------------------------------------------------------------------
%  ORCID
%------------------------------------------------------------------------------
% Creating some TikZ styles
\tikzset{
  nonterminal/.style = {rectangle
    , minimum size = 6mm
    , very thick
    , draw = black!
  }
}

% Create the ORCID logo
\definecolor{orcidlogocol}{HTML}{A6CE39}
\tikzset{
  orcidlogo/.pic={
    \fill[orcidlogocol]
      svg{M256,128c0,70.7-57.3,128-128,128C57.3,256,0,198.7,0,128C0,57.3,57.3
        ,0,128,0C198.7,0,256,57.3,256,128z};
    \fill[white]
      svg{M86.3,186.2H70.9V79.1h15.4v48.4V186.2z}
      svg{M108.9,79.1h41.6c39.6,0,57,28.3,57,53.6c0,27.5-21.5,53.6-56.8
        ,53.6h-41.8V79.1z M124.3,172.4h24.5c34.9,0,42.9-26.5
        ,42.9-39.7c0-21.5-13.7-39.7-43.7-39.7h-23.7V172.4z}
      svg{M88.7,56.8c0,5.5-4.5,10.1-10.1,10.1c-5.6,0-10.1-4.6-10.1-10.1c0-5.6
      ,4.5-10.1,10.1-10.1C84.2,46.7,88.7,51.3,88.7,56.8z};
  }
}

% Command to create the ORCID
\newcommand\orcid[1]{\href{https://orcid.org/#1}{\mbox{\scalerel*{
\begin{tikzpicture}[yscale=-1,transform shape]
  \pic{orcidlogo};
\end{tikzpicture}
}{|}}}}

%------------------------------------------------------------------------------
%	Titlepage: Header
%------------------------------------------------------------------------------
\title{Methods for Implementing Integrated Step-Selection Functions with
Incomplete Data}

% List of Authors
\author{
  David D. Hofmann\textsuperscript{1,2,\S} \orcid{0000-0003-3477-4365} \and
  Gabriele Cozzi\textsuperscript{1,2} \orcid{0000-0002-1744-1940} \and
  John Fieberg \textsuperscript{3} \orcid{0000-0002-3180-7021}
}

% Reduce spacing between authors
\makeatletter
\def\and{%
  \end{tabular}%
  \hskip -0.5em \@plus.17fil\relax
  \begin{tabular}[t]{c}}
\makeatother

% Current Date
% \date{\today}

% And here the masterpiece begins
\begin{document}

% Change page numbering
\pagenumbering{gobble}

% Create Titlepage
\maketitle

%------------------------------------------------------------------------------
%	Titlepage: Additional Info
%------------------------------------------------------------------------------
\begin{flushleft}

\vspace{0.5cm}

\textsuperscript{1} Department of Evolutionary Biology and Environmental
Studies, University of Zurich, Winterthurerstrasse 190, 8057 Zurich,
Switzerland.

\textsuperscript{2} Botswana Predator Conservation Program, Wild Entrust,
Private Bag 13, Maun, Botswana.

\textsuperscript{3} Department of Fisheries, Wildlife, and Conservation Biology,
University of Minnesota, St. Paul, MN, USA.

\textsuperscript{\S} Corresponding author: \href{mailto://david.hofmann2@uzh.ch}{david.hofmann2@uzh.ch}

\vspace{4cm}

\textbf{Running Title:} Step-Selection Analyses with Missing Data

\vspace{0.5cm}

\textbf{Keywords:} animal movement, gps data, imputation, incomplete data,
missing fixes, step-selection analyses, step-selection functions

\end{flushleft}

%------------------------------------------------------------------------------
%	Abstract
%------------------------------------------------------------------------------
\newpage
\begin{abstract}

Integrated step-selection analyses (iSSAs) are versatile and powerful frameworks
for studying habitat and movement preferences of tracked animals. iSSAs utilize
integrated step-selection functions (iSSFs) to model movements in discrete time,
and thus, require animal location data that are regularly spaced in time.
However, many real-world datasets are incomplete due to tracking devices failing
to locate an individual at one or more scheduled times, leading to slight
irregularities in the duration between consecutive animal locations. To address
this issue, researchers typically only consider bursts of regular data (i.e.,
sequences of locations that are equally spaced in time), thereby reducing the
number of observations used to model movement and habitat selection. We reassess
this practice and explore four alternative approaches that account for temporal
irregularity resulting from missing data. Using a simulation study, we compare
these alternatives to a baseline approach where temporal irregularity is ignored
and demonstrate the potential improvements in model performance that can be
gained by leveraging these additional data. We also showcase these benefits
using a case study on a spotted hyena (\textit{Crocuta crocuta}).

\end{abstract}

%------------------------------------------------------------------------------
%	Main Text
%------------------------------------------------------------------------------
% \newpage
%
% \onehalfspacing
% \tableofcontents
% \doublespacing

% Change page numbering
\newpage
\pagenumbering{arabic}

% Create linenumbers
\linenumbers

\section{Introduction}

Understanding how animals move across the landscape, what habitats they prefer,
and what resources they select are fundamental questions in movement ecology
\citep{Nathan.2008}. Thanks to recent advances in animal tracking
\citep{Cagnacci.2010, Williams.2019, Beardsworth.2022} and remote sensing
technologies \citep{Toth.2016, Rumiano.2020}, new opportunities and analytical
tools have emerged for studying how animals move and interact with their
environment \citep{Tomkiewicz.2010, Kays.2015, Nathan.2022}. Methods commonly
used to analyze animal movement data, including step-selection analyses
\citep{Fortin.2005, Thurfjell.2014, Fieberg.2021} and hidden Markov models
\citep{Michelot.2016}, require \textbf{animal locations} (terms in bold at first
occurrence are defined in \Cref{tab:Glossary}) that are collected at a constant
sampling frequency, leading to data that are equally spaced in time. Yet, it is
common to encounter missing locations in most telemetry data sets
\citep{Frair.2010, Hofman.2019, Vales.2022}, which introduces unwanted
irregularities in the duration between successive locations. Thus, there is a
need for analytical tools that enable the analysis of such data, while
mitigating potential biases arising from temporal irregularity introduced
through missing animal locations.

Step-selection analyses (SSAs) are widely used to study animals' movement
capacities and habitat-selection patterns \citep{Fortin.2005, Thurfjell.2014}.
Straight-line segments connecting consecutive animal locations, referred to as
\textbf{steps}, form the basic building blocks of the statistical likelihood in
SSAs. Specifically, SSAs model the probability $u$ of finding an individual at
location $s$ at time $t+1$, given the animal's past positions at time $t$ and
$t-1$, $s_t$ and $s_{t-1}$, respectively:

\begin{equation}
\label{EQ1}
u(s_{t+1}) = \frac{\phi(s_{t+1}, s_t, s_{t-1}; \gamma)w(x(s_{t+1}); \beta)}{\int_{s \in G}\phi(s_{t+1}, s_{t}, s_{t-1}; \gamma)w(x(s_{t+1}); \beta)ds}
\end{equation}

\noindent Here, the function $\phi$ represents an animal's \textbf{movement
kernel} which is usually expressed in terms of \textbf{step-length} and
\textbf{turning-angle} distributions, with $\gamma$ representing parameters in
these distributions. The function $w$ is the \textbf{habitat-selection function}
and reflects an animal's preferences $\beta$ for environmental characteristics
$x$ at location $s_{t+1}$. In most applications, $w$ is modeled as a log-linear
function of $x$, taking the form $w = exp(x^\top\beta)$. The integral in the
denominator of \ref{EQ1} ensures that $u$ is a proper probability distribution
(i.e., that it integrates to 1). Following \citet{Michelot.2024}, we call the
product $\phi \times w$ the step-selection function (SSF), as it highlights that
the probability of finding an animal at a certain location depends on both the
animal's movement kernel and its habitat-selection function.

Given a series of \textbf{observed steps}, finding the movement and
habitat-selection parameters that maximize the likelihood in \ref{EQ1} requires
approximating the integral in the denominator for each observed step. A variety
of numerical integration techniques can be used for this purpose
\citep{Michelot.2024}, but a common approach is to combine observed steps with
\textbf{random steps} generated by sampling step lengths and turning angles from
parametric distributions informed by the data \citep{Fortin.2005,
Thurfjell.2014}. Environmental conditions at observed steps are then contrasted
with environmental conditions at random steps in a (mixed effects) conditional
logistic regression framework \citep{Fortin.2005, Muff.2020}. To jointly
estimate parameters in $\phi$ and $w$, movement descriptors (e.g., step length
(sl), its natural logarithm (log(sl)), and the cosine of the turning angle
(cos(ta))) can be included in the conditional logistic regression model, and
their estimated coefficients can be used to update the initial (tentative)
step-length and turning-angle distributions \citep{Duchesne.2015, Avgar.2016,
Fieberg.2021}. The specific descriptors that need to be included depend on the
assumed step-length and turning-angle distributions \citep[for more details, see
Appendix C of][]{Fieberg.2021}. This approach to estimating parameters of the
SSF, termed \textit{integrated} SSA (or iSSA) by \citet{Avgar.2016}, is similar
to using importance sampling to approximate the integral in \ref{EQ1}
\citep{Michelot.2024} and is readily accessible through the \texttt{R}-package
\texttt{amt} \citep{Signer.2019}.

%In the SSF framework introduced by \citet{Fortin.2005}, the movement kernel
%$\phi$ was assumed to be known, which precluded joint estimation of $\phi$ and
%$w$ and has been shown to result in biased habitat-selection estimates
%\citep{Forester.2009}. A more robust alternative is to jointly estimate $\phi$
%and $w$ by fitting tentative distributions for step-lengths and turning-angles,
%and including descriptors of step-characteristics (i.e., step-lengths and
%turning-angles) as covariates in the statistical model \citep{Forester.2009,
%Avgar.2016}. In SSFs, the use of random steps serves to generate integration
%points that are needed to approximate the integral in \ref{EQ1} and, even
%though the true movement kernel $\phi$ is often not known a priori, it can be
%approximated using \textit{tentative} distributions \citep{Michelot.2024}. The
%use of parametric distributions serves as a means of conducting importance
%sampling and allows generating random steps that contribute better towards
%improving the approximation \citep{Michelot.2024}.

SSAs have proven extremely effective in numerous ecological studies
\citep{Thurfjell.2014}, providing insights into seasonal space use
\citep{Vales.2022, Enns.2023}, resource selection during distinct behavioral
phases \citep{Elliot.2014, Abrahms.2017, Cozzi.2018, Broekhuis.2019}, and the
effects of landscape familiarity or memory on animal movements \citep{Kim.2023}.
A model parametrized using iSSA resembles a fully mechanistic movement model
that can be used to simulate space use under novel conditions \citep{Avgar.2016,
Signer.2017, Hofmann.2023, Signer.2023}. This characteristic has made iSSAs a
useful tool for quantifying landscape resistance and identifying movement
corridors \citep{Zeller.2020, Buchholtz.2020, Hofmann.2021, Hofmann.2023}.

% Glossary of terms
\begin{table}

\caption{\label{tab:Glossary}Glossary of terms. Terms in the glossary are printed in bold at first occurrence in the main text. Definitions are always given in the context of step selection functions (SSFs).}
\centering
\resizebox{\linewidth}{!}{
\begin{tabular}[t]{>{}l>{\raggedright\arraybackslash}p{10 cm}}
\toprule
Term & Definition\\
\midrule
\textbf{\cellcolor{gray!6}{Animal locations}} & \cellcolor{gray!6}{A series of telemetry data points that include date, time, longitude, and latitude information, describing when and where an animal was observed or recorded.}\\
\textbf{Step} & A straight line connecting two consecutive locations.\\
\textbf{\cellcolor{gray!6}{Observed step}} & \cellcolor{gray!6}{A step that connects two observed animal locations.}\\
\textbf{Random step} & A step that connects an observed animal location with a random location. Random locations are typically generated by combining an observed animal location with random step lengths and turning angles.\\
\textbf{\cellcolor{gray!6}{Step length}} & \cellcolor{gray!6}{The Euclidean distance of a step.}\\
\addlinespace
\textbf{Turning angle} & A measure of the change in direction between two consecutive steps.\\
\textbf{\cellcolor{gray!6}{Habitat-selection function}} & \cellcolor{gray!6}{A probabilistic description of an animal's habitat preferences. It describes how an animal selects habitat when not constrained by its movement capacity. Also referred to as movement-free habitat-selection function.}\\
\textbf{Movement kernel} & A probabilistic description of an animal's movement capacity. It describes how an animal would move when not constrained by habitat selection. Also referred to as selection-free movement kernel.\\
\textbf{\cellcolor{gray!6}{Trajectory}} & \cellcolor{gray!6}{A sequence of animal locations collected on the same individual.}\\
\textbf{Step duration} & The time interval associated with a particular step, i.e., the time elapsed between two consecutive animal locations.\\
\addlinespace
\textbf{\cellcolor{gray!6}{Regular animal locations}} & \cellcolor{gray!6}{A series of animal locations that have been obtained at regularly spaced time intervals, such as every hour.}\\
\textbf{Irregular animal locations} & A series of animal locations collected at irregular time intervals.\\
\textbf{\cellcolor{gray!6}{Regular step durations}} & \cellcolor{gray!6}{Step durations that occur when animal locations are successfully collected at regular time intervals.}\\
\textbf{Irregular step durations} & Step durations that occur when animal locations are not successfully collected at regular time intervals.\\
\textbf{\cellcolor{gray!6}{Missingness}} & \cellcolor{gray!6}{The fraction of animal locations that should have been collected but, for some reason, were not. For example, if only eight out of ten expected animal locations were successfully collected, the missingness would be 0.2 (i.e., 20\%).}\\
\addlinespace
\textbf{Forgiveness} & The maximum step duration, measured in multiples of the regular step duration, a modeler is willing to include in the step-selection analysis. A modeler with a forgiveness of one, for instance, only considers regular steps, while a modeler with a forgiveness of two would consider irregular steps up to twice the regular step duration.\\
\textbf{\cellcolor{gray!6}{Burst}} & \cellcolor{gray!6}{A sequence of consecutive animal locations equally spaced in time and with steps were the step duration does not exceed the forgiveness.}\\
\textbf{Valid step} & A step for which a step length, turning angle, and step duration can be computed, and for which the step duration does not exceed the forgiveness of the modeler. These steps can be used for step-selection analysis.\\
\bottomrule
\end{tabular}}
\end{table}


A key assumption when conducting an iSSA is that the data have been collected at
a constant sampling frequency, thus producing \textbf{trajectories} with
\textbf{regular step durations} ($\Delta t$; \citealp{Fortin.2005,
Thurfjell.2014}). Here, we refer to such data as \textbf{regular animal
locations}, and without loss of generality, we assume the regular step duration
to be one (i.e., $\Delta t = 1$). Regular step durations ensure that
step lengths and turning angles are independent of the step duration, and
therefore, steps can be pooled when estimating movement parameters.
Since animal locations are usually obtained using automated tracking devices,
such as GPS collars programmed to record data at regular intervals, satisfying
this assumption seems straightforward. In reality, however, device limitations
often imply that some of the aspired datapoints fail to be collected, thus
introducing \textbf{missingness} and confronting researchers with
\textbf{irregular animal locations} and \textbf{irregular step durations}
\citep{Frair.2010}. In a comprehensive study, \citet{Hofman.2019} showed that
across 3,000 GPS devices and 160 species, the average success rate of obtaining
a scheduled animal location was 78\% (implying a missingess of 22\%), thus
highlighting that irregular animal locations are a frequent phenomenon in
ecological studies.

It is generally recommended that, in the case of such irregular data,
researchers should only retain \textbf{bursts} of steps with regular
step durations (possibly with some tolerance) and discard the rest
\citep{Thurfjell.2014}. We will refer to this modeling approach as having a
\textbf{forgiveness} level of one, indicating that only steps with
step durations $\Delta t = 1$ are retained for further analysis. In \texttt{R},
the \texttt{amt} package provides the function \texttt{track\_resample}
specifically for identifying bursts of steps with regular step durations
\citep{Signer.2019}. The main drawback of this approach is that it may result in
a substantial amount of data being discarded (\Cref{Overview} and
\Cref{NumberOfSteps}). For instance, consider a hypothetical trajectory in which
location 4 is missing (\Cref{Overview}a). The absence of this location prevents
the computation of a step between locations 3 and 4, as well as between
locations 4 and 5. Furthermore, without these steps, it becomes impossible to
compute a turning angle for the step between locations 5 and 6. Consequently,
the lack of a single location reduces the effective sample size, which is the
number of \textbf{valid steps}, by three. Assuming that animal locations are
missing at random, a missingness level of 25\% causes the number of valid steps
to drop by $58\%$ (\Cref{NumberOfSteps}). A modeler willing to increase their
level of forgiveness to two (i.e., allowing for inclusion of steps with $\Delta
t \leq 2$) would be able to increase the number of valid steps by $57\%$
(\Cref{Overview} and \Cref{NumberOfSteps}), therefore achieving a substantial
gain in effective sample size. The ability to capitalize on irregular data is
likely to be particularly important for applications where data are already
limited, such as, for instance, when modeling dispersing individuals
\citep{Rudnick.2012, Fattebert.2015, Cozzi.2020}. However, increasing the
forgiveness also implies that step durations of the retained steps become
irregular, thus necessitating appropriate tools to account for such
irregularity.

\begin{figure}
  \begin{center}
  \includegraphics[width = \textwidth]{99_Overview.png}
  \caption{(\textbf{a}) Demonstration of how missingness affects the number of
  valid steps that can be used for step-selection analyses under different
  levels of forgiveness. The upper panel depicts a trajectory with zero missing
  locations. That is, all aspired locations were successfully collected on a
  regular interval (yielding a regular step duration of $\Delta t = 1$). This
  trajectory produces four valid steps that can be included in the iSSF model
  and one invalid step that has to be omitted because it has no turning angle
  associated with it. In the central panel, animal location 4 was not obtained,
  introducing a missingness of 16.7\%. If the modeler has a forgiveness of one,
  only a single step can be included for further analysis, as all other steps
  are invalid (either because no turning angle can be computed or because
  step durations exceed the forgiveness). If, however, the modeler exhibited a
  forgiveness of two, such as in the lower panel, a total of three steps could
  be obtained for further analysis. (\textbf{b}) Conceptual illustration of how
  increasing the forgiveness allows one to retain additional steps that can be
  used for step-selection analysis. The sequence of dots resembles the sequence
  of locations that were scheduled to be collected (e.g. using a GPS device),
  with the lines representing hypothetical steps. Because not all locations were
  successfully obtained (gray dots), there is missingness. Depending on the
  forgiveness level, already a single missing location enforces the introduction
  of a new burst, which leads to the loss of several steps. In addition, some of
  the remaining steps are invalid (dotted) because they are lacking a
  turning angle. By increasing the forgiveness, a modeler is willing to retain
  steps that exceed the regular step duration by a certain threshold, which
  enables them to obtain longer bursts and increase the number of steps that can
  be used for further analysis. In the figure, forgiveness increases from left
  to right.}
  \label{Overview}
  \end{center}
\end{figure}

\begin{figure}
  \begin{center}
  \includegraphics[width = \textwidth]{99_NumberOfSteps.png}
  \caption{Illustration of how missingness in animal location data reduces the
  number of valid steps that can be used in step-selection analyses (left panel)
  and how increasing forgiveness helps to retain additional steps that are
  otherwise omitted (right panel). At a missingness of 0, 998 valid steps can be
  obtained from the total of 1,000 animal locations. At higher missingness,
  step durations become irregular, which means that the number of valid steps
  decreases substantially. However, if the modeler is willing to increase their
  forgiveness, additional steps can be gained. The right panel shows the number
  of valid steps that is gained when increasing the forgiveness from 1 to
  2, 3, 4, and 5, respectively. Ribbons indicate the 95\%-percentile intervals
  derived from 1,000 replicates.}
  \label{NumberOfSteps}
  \end{center}
\end{figure}

Various methods have been employed in the past to address temporal
irregularities in animal location data. These may serve as valuable starting
points for developing approaches that enable the integration of irregular data
with iSSFs.

\begin{enumerate}

  \item \textit{Imputation}: An intuitive solution is provided by
  \citet{McClintock.2017}, who suggests fitting a continuous-time correlated
  random walk movement model \citep{Johnson.2008} to the collected data and to
  use the fitted model to impute missing fixes. By imputing missing locations,
  the analysed trajectories become entirely regular again and can be analysed
  using traditional techniques. This approach, which we coin the
  \textit{imputation} approach, is readily available through the
  \texttt{R}-package \texttt{crawl} \citep{Johnson.2022}, yet has only been
  tested for use with hidden Markov movement models and not with iSSFs
  \citep{McClintock.2017}.

  \item \textit{Na\"ive}: Another approach is outlined by \citet{Munden.2021},
  who introduced time-varying iSSFs. In this framework, a change-point detection
  algorithm is applied to the series of observed animal locations to identify
  distinct decision points where the animal turns \citep{Potts.2018,
  Munden.2021}. Steps are then created to represent straight-line movements
  in-between these decision points, but because decision points are not
  regularly spaced in time, the resulting step durations are irregular. Thus,
  step durations are treated as random variables, and, instead of generating
  random steps by sampling step lengths and turning angles, the authors generate
  random steps by sampling step durations, step speeds, and turning angles. The
  underlying assumption is that step lengths scale linearly with step duration,
  and can therefore be meaningfully represented by combinations of step speeds
  and step durations. Although this approach was developed with
  ultra-high-frequency data in mind, we might \textit{na\"ively} apply it more
  broadly to the case of missing data if we believe the presumed
  linear relationship to hold true (this assumption may be reasonable when
  step durations are short, but it is unlikely to hold for longer step durations
  since an animal's path will deviate from a straight line between successive
  locations). Hence, we propose, with our \textit{na\"ive} approach, to scale the
  generated random steps by the observed step duration.

  \item \textit{Dynamic+Model}: Instead of generating random steps by sampling
  step lengths and turning angles from distributions fitted to a single
  step duration, one may choose to fit separate distributions to steps of
  different durations, thus acknowledging potentially non-linear relationships
  between step duration and step lengths or turning angles. Because random steps
  in this approach are sampled using different \textit{tentative} distributions,
  it is necessary to include interactions between step durations and other
  step descriptors (e.g., sl, log(sl), and cos(ta)) in the conditional logistic
  regression model to allow updating the distributions to the different
  step durations. We therefore refer to this approach as the
  \textit{dynamic+model} approach, highlighting that step-length and
  turning-angle distributions are dynamically adjusted to observed
  step durations and that the step duration is included as a modifier of the
  coefficients of step descriptors in the regression model.

  \item \textit{Multistep}: Finally, we propose a \textit{multistep} approach,
  where random steps of varying step durations are generated by stitching
  together sequences of random steps from the regular step duration. For
  example, one can generate a random step of duration $\Delta t = 2$ by
  combining two random steps of step duration $\Delta t = 1$.

\end{enumerate}

\noindent Our goal with this article is to reassess the practice of discarding
irregular animal locations in iSSFs and to investigate whether retaining
irregular data could, in fact, serve to improve model performance. Our
hypothesis is that even irregular data contains valuable information on habitat
and movement preferences that could be leveraged if appropriate methods are
applied. To test this notion, we conducted a comprehensive simulation study
where we simulated regular animal location data with known movement and habitat
parameters. We then introduced varying levels of missingness and applied iSSFs
to estimate simulation parameters. Specifically, we employed the four
alternative iSSF approaches outlined above and compared them to the traditional
approach of including only bursts of regular data and to an uncorrected approach
that simply ignored irregular step durations when using a forgiveness level $>$
1. To examine the impact of different landscape configurations on derived
estimates, we ran our simulations for different levels of spatial
autocorrelation. The use of simulations instead of real data had the benefit
that underlying parameters of the movement kernel and habitat-selection function
were known, which allowed us to assess the reliability of different methods in
retrieving true simulation parameters under different conditions (sensu
\citealp{Kery.2016}). We then compare the traditional approach (using only
bursts of steps with regular step durations) to the best-performing approach
(with irregular data) using GPS locations collected on a spotted hyena
(\textit{Crocuta crocuta}).

We anticipated that increasing forgiveness without adjusting for the introduced
irregularity would entail a bias-variance trade-off. Specifically, we
anticipated that increasing forgiveness would allow improving estimator
precision, but at the cost of introducing bias due to failing to account for
irregular sampling intervals. We expected this bias to be particularly
pronounced at high levels of missingness. Furthermore, we hypothesized that
accounting for irregularity in the \textit{na\"ive}, \textit{dynamic+model}, and
\textit{multistep} approaches would improve model accuracy, while alleviating
potential bias, thus providing an effective means of incorporating additional
data. Because the \textit{imputation} approach relied on an intermediate
movement model to predict missing animal locations, we had no prior expectations
for how well it would perform.

\section{Methods}

We implemented the simulation study in the programming language \texttt{R}
version 4.3.2 \citep{RCoreTeam.2023} and achieved parallelization of
simulation-runs using the \texttt{R}-package \texttt{pbmcapply}
\citep{Kuang.2022}. We generated figures using the \texttt{ggplot2}
\citep{Wickham.2023}, \texttt{ggpubr} \citep{Kassambara.2023}, and
\texttt{ggh4x} \citep{Brand.2023} \texttt{R}-packages. We manipulated raster
data and computed spatial distances using the \texttt{R}-package \texttt{raster}
\citep{Hijmans.2023}. An overview of the simulation design and the different
iSSF approaches is presented in \Cref{Design} and all codes to reproduce this
study are available through an online repository \citep{Hofmann.2023a}.

\begin{figure}
  \begin{center}
  \includegraphics[width = \textwidth]{99_Design.png}
  \caption{Illustration of the study design. We varied the autocorrelation range
  when simulating spatial covariates from 5 to 20 and tested for different
  missingness scenarios (ranging from 0\% to 50\% missing locations). To
  investigate how increasing forgiveness (i.e., the willingness to include steps
  with duration above the regular step duration) influenced model results, we
  varied its value from 1 (regular step selection) to 5 (considering steps that
  are five times the regular step duration). Finally, we tested five different
  methods to account for potential biases introduced by including irregular
  steps. This gave 3 x 6 x 5 x 5 = 450 combinations, each of which we replicated
  100 times. We assumed step lengths (sl) to follow a gamma distribution,
  whereas turning angles (ta) followed a von Mises distribution.}
  \label{Design}
  \end{center}
\end{figure}

\subsection{Landscape Simulation}

We simulated a virtual landscape comprising two continuous and one categorical
(binary) spatial layers, each with a resolution of 300 x 300 pixels
(\Cref{Covariates}) and spanning across x- and y-coordinates from 0 to 300. The
first layer, \textsf{Dist} (continuous), quantified the distance to the center
of the virtual landscape ($x = 150$, $y = 150$), and can be understood as a
point of attraction, such as, for instance, the center of an animal's
home-range. The second layer, \textsf{Elev} (continuous), resembled an elevation
layer and was simulated by sampling random pixel-values from a normal
distribution. To achieve spatial autocorrelation, we applied a circular moving
window with radius $r$ within which we tallied pixel-values. We varied $r$ from
5, to 10, to 20, depending on the simulated level of autocorrelation (Figure
S1). The third layer, \textsf{Forest} (categorical), represented areas covered
by woodland and was simulated similarly to the \textsf{Elev} layer, but we
binarized the layer by setting all simulated values above the 50\% quantile to
forest and all other values to non-forest (our reference class). We normalized
values of all simulated layers to a range between zero and one and replicated
the simulation of each layer 100 times per autocorrelation scenario, thus
producing 300 unique landscape configurations.

\begin{figure}
  \begin{center}
  \includegraphics[width = \textwidth]{99_Covariates.png}
  \caption{Example of a landscape configuration across which we simulated
  movement trajectories. All simulated layers had a resolution of 300 x 300
  pixels. The distance layer indicated the distance to the center of the
  landscape and served to simulate attraction. The elevation and forest layers
  were simulated by sampling pixel-values from a normal distribution and
  applying a moving window to achieve spatial autocorrelation. Simulated
  individuals were initiated within the white dashed rectangle, which ensured
  that they would not be released directly at a map border. Simulated
  individuals were attracted to the landscape's center, preferred elevated
  areas, and avoided areas covered by forest. The black line shows the simulated
  trajectory associated with the visualized landscape configuration (cfr.
  \Cref{MovementSimulation}).}
  \label{Covariates}
  \end{center}
\end{figure}

\subsection{Movement Simulation}
\label{MovementSimulation}
To simulate movement across the virtual landscape, we employed the iSSF
simulation algorithm developed by \citet{Signer.2017} and applied in
\citet{Hofmann.2023}. This procedure consists of a sequence of five steps that
are repeated $n$ times to generate a movement trajectory. In step one, we
generated a random starting location by sampling random x- and y-coordinates on
the simulated landscape. To prevent starting points near map borders, we
restricted sampled locations to x- and y-coordinates between 50 and 250 (white
dotted rectangle in \Cref{Covariates}). In step two, we generated a set of 1,000
random steps originating from the current location, by sampling turning angles
from a von Mises distribution with concentration parameter $\kappa = 0.5$ and
step lengths from a gamma distribution with shape parameter $k = 3$ and scale
parameter $\theta = 1$. In step three, we extracted covariate values at the end
of each random step from the underlying covariate layers. In step four, we
assigned to each step $j$ a probability $\pi_j$ of being selected using the
equation below:

\begin{equation}
\label{EQ2}
  \pi_j =
  \frac{exp(\beta^\top x_j)}{\sum_{i = 1}^J exp(\beta^\top x_i)} %=
  % \frac{
  %   exp(\beta_{dist} dist_j + \beta_{elev} elev_j + \beta_{forest} forest_j)
  % }{
  %   \sum_{i = 1}^J exp(\beta_{dist} dist_i + \beta_{elev} elev_i + \beta_{forest} forest_i)
  % }
\end{equation}

\noindent Here, $\beta$ represents the vector of habitat-selection parameters
and $x_j$ the covariate value at the end of the $j$\textsuperscript{th} step.
The probability of a step being selected thus depended on its associated
covariate values, the covariate values of all other random steps, and the
simulated preferences $\beta$. We defined the habitat-selection parameters as
$\beta_{dist} = -20$, $\beta_{elev} = 0.5$, and $\beta_{forest} = -1$. That is,
simulated individuals were attracted to the landscape's center, preferred
elevated areas, and avoided areas covered by forest. In step five, we sampled
one of the random steps based on predicted probabilities and computed the
simulated individual's new position. We then repeated steps two through five
until the trajectory comprised a total of 1,000 steps. Each simulated step was
assumed to have a step duration of exactly $\Delta t = 1$. We repeated the
simulation for each of the 300 simulated landscapes, producing 300 unique
trajectories (example trajectory presented in \Cref{Covariates}).

\subsection{Data Rarefication}

To simulate missingness, we rarefied the trajectories by randomly removing a
fixed fraction of animal locations. To assess the impact of different degrees of
missingness, we varied the fraction of removed data from 0\% (complete dataset)
to 50\% in increments of 10\%. The random removal of animal locations introduced
temporal irregularity, such that the resulting step durations differed depending
on the time elapsed between remaining fixes. We replicated the rarefication of
each trajectory 100 times.

\subsection{Computing Bursts}

We used the rarefied data to compute bursts consisting of a sequence of
animal locations with step durations that did not exceed the forgiveness value.
To test how different levels of forgiveness impacted our results, we varied
forgiveness from 1 (maximum allowed step duration was $\Delta t = 1$) to 5
(maximum allowed step duration was $\Delta t = 5$). As an example, if the
forgiveness was 1, any step with step duration $\Delta t > 1$ resulted in a new
burst. If the forgiveness was 2, in contrast, step durations of up to $\Delta t
= 2$ were allowed before a new burst was introduced (\Cref{Overview}b). Within
each burst, we calculated step lengths and turning angles. However, due to the
grouping of steps into bursts, the orientation of the first step within each
burst relative to the previous step could not be determined. As a result, this
step always lacked a turning angle and was considered invalid
(\Cref{Overview}b).

\subsection{Fitting Distributions}

Based on the steps retained within bursts, we parametrized tentative step-length
and turning-angle distributions. Specifically, we used the \texttt{fit\_distr}
function from the \texttt{amt} package \citep{Signer.2019}, which is a wrapper
function for the \texttt{fitdist} function from the \texttt{fitdistrplus}
package \citep{Delignette-Muller.2015}, and fitted a gamma distribution to
step lengths and a von Mises distribution to turning angles. Notably, we
employed two different fitting procedures:

\begin{enumerate}

  \item \textit{Regular Distributions:} In this procedure, we fitted parametric
  distributions considering only step lengths and turning angles from steps that
  exhibited a step duration of $\Delta t = 1$ (i.e., the regular step duration).
  Any steps with irregular step durations ($\Delta t > 1$) were discarded and
  did not affect distributional parameter estimates. This represents the
  traditional procedure in iSSFs where only regular bursts of animal locations
  are considered when estimating tentative movement parameters.

  \item \textit{Dynamic Distributions:} In this procedure, we fitted separate
  parametric distributions to step lengths and turning angles from steps of
  different step durations. That is, we parametrized separate turning-angle and
  step-length distributions representative of steps with durations of $\Delta t
  = 1, 2, 3, 4$ and $5$ (which corresponds to the maximum forgiveness level we
  tested for). Some step durations only rarely occurred at low levels of
  missingness, thus complicating parametrization of the associated
  distributions. To facilitate estimation of dynamic distribution parameters
  across all $\Delta t$ (Figure S2), we resampled data to different step
  durations using the \texttt{track\_resample} function from the \texttt{amt}
  package \citep{Signer.2019} before fitting tentative parameters. This ensured
  a sufficient number of steps for each step duration to estimate associated
  parameters. An alternative approach would be to increase missingness in the
  data even further, thus introducing a larger number of longer step durations.

\end{enumerate}

\subsection{Step-Selection Functions}

We implemented a baseline \textit{uncorrected} approach and four alternative
iSSF approaches that mainly differed in the way in which random steps were
generated, but sometimes also in the model call that was used to estimate
parameters (\Cref{Design}). In the \textit{uncorrected} approach, we treated
data as if they were regular, ignoring potential issues arising from having
variable step durations. When forgiveness was one, this approach corresponded
to the traditional iSSF approach. All other approaches were targeted towards
reducing potential biases arising from the inclusion of steps with irregular
step durations. Irrespective of the approach employed, we paired each observed
step with a total of 200 random steps:

\begin{itemize}

  \item \textit{Uncorrected}: In the uncorrected approach, we generated random
  steps by sampling step lengths and turning angles from statistical
  distributions fitted to steps with step durations of $\Delta t = 1$,
  regardless of the forgiveness value or observed step durations. Thus, this
  approach ignored any potential effect of step duration when generating random
  step lengths and turning angles.

  \item \textit{Imputed}: In this approach, we sampled step lengths and turning
  angles from statistical distributions fitted to observed steps with a step
  duration of $\Delta t = 1$. However, prior to generating random steps, we
  imputed missing fixes using predictions from a simple movement model.
  Specifically, we fitted a single-state movement model \citep{Johnson.2008} to
  the simulated trajectories and used the parametrized model to predict
  coordinates for all missing animal locations. For this, we used the functions
  \texttt{crwFit} and \texttt{crwPredict} from the \texttt{crawl}
  \texttt{R}-package \citep{Johnson.2022}. Although the \texttt{crwFit} function
  provides capabilities to incorporate location measurement error, we assumed
  animal locations were measured without error. The imputation resulted in a
  complete dataset without any missing animal locations, such that each
  trajectory consisted of a single continuous burst of locations equally spaced
  in time. As such, the imputation approach is not affected by the forgiveness
  level.

  \item \textit{Na\"ive}: In the \textit{na\"ive} approach, we again sampled
  step lengths and turning angles from regular distributions fitted to steps
  with step durations of $\Delta t = 1$. However, we linearly scaled the sampled
  step lengths depending on the step durations of the observed steps. For
  instance, we doubled the sampled step lengths for any observed step with a
  step duration of $\Delta t = 2$. This approach na\"ively assumed that
  step lengths scale linearly with step durations, which is unlikely to be true,
  as most animals don't move in straight lines between successive observations.
  Furthermore, the linear approximation is likely to get worse as step duration
  increases (i.e., as the forgiveness value increases). Since it is not clear
  how turning angles should scale with step duration, we did not adjust the
  sampled turning angles.

  \item \textit{Dynamic+Model}: In the \textit{dynamic+model} approach, we
  sampled step lengths and turning angles from dynamic distributions that were
  fit to different step durations. That is, for observed steps with
  step duration of $\Delta t = 2$, we sampled step lengths and turning angles
  from distributions fit to observed steps with $\Delta t = 2$. We then included
  interactions between the step duration and other step descriptors (e.g., sl,
  log(sl), cos(ta)), allowing us to update movement parameters for each
  step duration separately. To avoid numerical instabilities with the
  conditional logistic regression model, we only included steps with durations
  $\Delta t > 1$ if the respective duration was represented at least 5 times in
  the rarefied dataset.

  \item \textit{Multistep}: In the \textit{multistep} approach, we sampled
  step lengths and turning angles from statistical distributions fitted to
  observed steps with step durations of $\Delta t = 1$. We then generated a
  sequence of random steps such that their combined step duration equaled the
  step duration of each observed step.  For instance, for an observed step with
  step duration of $\Delta t = 2$, we generated sets of two random steps, which
  we then concatenated into a ``random path''. The paths were then simplified to
  straight lines connecting the first and last coordinate of each path, which
  represented the final random step.

\end{itemize}

\noindent Together, an observed step and its 200 associated random steps formed
a \textit{stratum} that received a unique ID. At the end of each step, we
extracted covariate values from the underlying covariate layers.

\subsection{Conditional Logistic Regression Model}

We estimated movement and habitat-selection parameters for the simulation
scenarios presented in \Cref{Design} using conditional logistic regression,
implemented using the \texttt{clogit} function in the \texttt{R}-package
\texttt{survival} \citep{Therneau.2023}. We defined a binary response variable
(\textsf{observed}) indicating if a step was an observed (scored 1) or a random
step (scored 0) and used the step's ID as a stratification variable. We included
habitat covariates (dist, elev, forest) and step descriptors (sl, log(sl),
cos(ta)) as predictors in the regression model. For the \textit{dynamic+model}
approach, we also included interactions between the step duration, coded as a
factor, and step descriptors. To update tentative movement parameters (denoted
by the subscript $_0$) and obtain the selection-free movement kernel (denoted by
the \string^ symbol), we employed the formulas provided in \citep{Avgar.2016,
Fieberg.2021}. Specifically, we updated the shape ($\hat{k}$) and scale
($\hat{\theta}$) parameters of the step-length distribution (gamma) using:

$$
\hat{k} = k_0 + \beta_{log(sl)}
$$

$$
\hat{\theta} = \frac{1}{\frac{1}{\theta_0} - \beta_{sl}}
$$

\noindent We updated the concentration parameter ($\hat{\kappa}$) of the
turning-angle distribution (von Mises) using:
$$
\hat{\kappa} = \kappa_0 + \beta_{cos(ta)}
$$

\noindent We kept track of the  estimates of the updated movement  ($\hat{k}$,
$\hat{\theta}$, and $\hat{\kappa}$) parameters and  the habitat-selection
($\hat{\beta}_{dist}$, $\hat{\beta}_{elev}$, $\hat{\beta}_{forest}$) parameters,
and compared them to the true simulation parameters. We also quantified model
accuracy via the root-mean-square error (RMSE).

\section{Results}

Results were qualitatively similar for all three landscape autocorrelation
scenarios and for different combinations of missingness and forgiveness (Figure
S4). Here, we report on results for a landscape with autocorrelation of 20,
while either holding constant missingness at a conservative 20\%
(\Cref{ModelComparisonForgiveness}) or the forgiveness level at two
(\Cref{ModelComparisonMissingness}) (results for all other combinations are
summarized in Figure S4). The \textit{imputation} approach resulted in biased
estimators of $\beta_{dist}$ and $\beta_{forest}$, whereas all other approaches
were able to recover the parameters of the habitat-selection function with
minimal bias (\Cref{ModelComparisonForgiveness}). Note, the \textit{imputation}
approach always starts with a full trajectory and is therefore unaffected by the
forgiveness level. For all other methods, increasing the forgiveness from 1 to 5
improved the precision of the estimators of habitat-selection parameters without
introducing noticeable bias, with the biggest gains in precision and RMSE
occurring when moving from a forgiveness of one to a forgiveness of two
(\Cref{ModelComparisonForgiveness}). This highlights the potential benefits of
leveraging additional data compared to the traditional approach, which uses only
bursts of regular data (represented by the \textit{uncorrected} approach and
forgiveness = 1).

The \textit{uncorrected}, \textit{na\"ive}, and \textit{imputation} approaches
resulted in biased estimators of the parameters in the movement kernel,
particularly for high values of forgiveness (\Cref{ModelComparisonForgiveness}a)
and high levels of missingness (\Cref{ModelComparisonMissingness}a). The
\textit{imputation} approach appeared to perform particularly poorly at
estimating the concentration parameter of the turning-angle distribution
(\Cref{ModelComparisonMissingness}). The \textit{multistep} and
\textit{dynamic+model} approaches resulted in unbiased estimators of parameters
in the step-length distribution, but estimators of the concentration parameter
exhibited a slight bias. This bias was, however, much smaller than we observed
with the other approaches we considered. Increasing missingness negatively
influenced the precision and accuracy of estimates, yet its impact could be
dampened using the \textit{dynamic+model} and \textit{multistep} approaches
(\Cref{ModelComparisonMissingness}b).

\begin{figure}
  \begin{center}
  \includegraphics[width = \textwidth]{99_ModelComparisonForgiveness.png}
  \caption{(\textbf{a}) Parameter estimates and (\textbf{b}) Root mean-square
  error (RMSE) with regard to the movement kernel and habitat-selection function
  as a function of forgiveness. Results are shown for the scenario with
  landscape autocorrelation of 20 and missingness of 20\%. The movement kernel
  comprised of a gamma distribution with shape parameter $k$ and scale parameter
  $\theta$ governing the step-length distribution and a von Mises distribution
  with concentration parameter $\kappa$ governing the turning-angle
  distribution. Habitat-selection was based on three covariates, namely a
  \textsf{Distance}, \textsf{Elevation}, and a \textsf{Forest} layer. Estimates
  are shown for the five different approaches we tested for. The uncorrected
  approach ignored the fact that higher forgiveness implied temporal
  irregularity in the data, while all other approaches attempted to correct for
  the potential biases introduced by temporal irregularity. Note, the imputation
  approach is not affected by the forgiveness level, since it always starts with
  a full trajectory.}
  \label{ModelComparisonForgiveness}
  \end{center}
\end{figure}

\begin{figure}
  \begin{center}
  \includegraphics[width = \textwidth]{99_ModelComparisonMissingness.png}
  \caption{(\textbf{a}) Parameter estimates and (\textbf{b}) Root mean-square
  error (RMSE) with regard to the movement kernel and habitat-selection function
  as a function of missingness. Results are shown for the scenario with
  landscape autocorrelation of 20 and forgiveness of 2. The movement kernel
  comprised of a gamma distribution with shape parameter $k$ and scale parameter
  $\theta$ governing the step-length distribution and a von Mises distribution
  with concentration parameter $\kappa$ governing the turning-angle
  distribution. Habitat-selection was based on three covariates, namely a
  \textsf{Distance}, \textsf{Elevation}, and a \textsf{Forest} layer. Estimates
  are shown for the five different approaches we tested for. The uncorrected
  approach ignored the fact that higher forgiveness implied temporal
  irregularity in the data, while all other approaches attempted to correct for
  the potential biases introduced by temporal irregularity. Note, the imputation
  approach is not affected by the forgiveness level, since it always starts with
  a full trajectory.}
  \label{ModelComparisonMissingness}
  \end{center}
\end{figure}

\section{Case Study}

To showcase the applicability of the \textit{dynamic+model} approach, which
appeared to perform best with simulated data, we conducted a case study with
real GPS data obtained on ``Apollo'', a spotted hyena (\textit{Crocuta crocuta})
inhabiting the Okavango Delta ecosystem in northern Botswana. Apollo's data were
collected between 2007 and 2011 using GPS radio collars (\textit{GPS Plus};
Vectronic Aerospace GmbH, Berlin, Germany) and comprised 9,316 GPS locations
(details in \citealp{Cozzi.2013a} and \citealp{Cozzi.2015}). Because hyenas are
nocturnal \citep{Cozzi.2012}, GPS collars were set to record data at two-hourly
intervals between 18:00 and 06:00 o'clock, and to only record a single location
at noon. For simplicity, we only considered nightly bursts and removed all
locations obtained at noon. Missingness in this dataset was low ($<$ 10 \%,
\citealp{Cozzi.2015}) and to better showcase the usefulness of the
\textit{dynamic+model} approach, we thinned the data by randomly removing 25\%
of the obtained locations. As spatial covariate layers, we used \textsf{Water},
\textsf{DistanceToWater} and \textsf{Trees} (Figure S5). \textsf{Water} was a
binary variable representing major rivers and areas inundated by floodwater,
whereas \textsf{DistanceToWater} was a continuous variable indicating the
distance (in meters) to the nearest pixel categorized as water. \textsf{Trees}
was a continuous variable indicating the percent tree cover in each pixel. We
resampled all layers to a common resolution of 250 m by 250 m and merged them
into a single raster-stack (Figure S5). The derivation of each covariate layer
is described in detail in \citealp{Hofmann.2021}. We dynamically fitted
step-length and turning-angle distributions to steps with step durations of 2,
4, and 6 hours, respectively, assuming a gamma distribution for step lengths and
a von Mises distribution for turning angles. Instead of resampling the observed
track to different step durations when fitting dynamic distributions (like we
did in the simulation study), we introduced a larger amount of steps with step
durations longer than two hours by thinning the data again (by another 10\%).
The benefit of this approach was that steps with irregular step durations
occurred more randomly and were not limited to the hours specified by the
resampling algorithm. Finally, we used iSSFs with the \textit{dynamic+model}
approach to estimate the habitat-selection function and movement kernel of
Apollo. For this, we considered three cases:

\begin{itemize}
  \item \textit{F1}: We assumed a forgiveness of one (i.e., only steps with a
  regular step duration of 2 hours), which is akin to conducting a traditional
  iSSA.

  \item \textit{F3-S}: We assumed a forgiveness of three (i.e., considered steps
  with step durations of up to three times the regular step duration) and
  included interactions between the step duration ($\Delta t$) and step
  descriptors (sl, log(sl), and cos(ta)) in the regression model.

  \item \textit{F3-SH}: We assumed a forgiveness of three (i.e., considered
  steps with step durations of up to three times the regular step duration) and
  included interactions between the step duration ($\Delta t$) and step
  descriptors (sl, log(sl), and cos(ta)), as well as between the step duration
  and habitat covariates in the regression model.

\end{itemize}

\noindent Notably, we included F3-SH to investigate if including interactions
between the step duration and habitat-covariates would provide insights into
scale-dependent habitat selection. In all cases, we generated 200 random steps
and extracted spatial covariates at the end of observed and random steps. We
then fit the three models using the conditional logistic regression framework as
implemented in the \texttt{survival} \texttt{R}-package \citep{Therneau.2023}.
Lastly, we computed updated movement parameters for a regular step duration of 2
hours.

Results from the iSSF models show that increasing the level of forgiveness led
to improvements in estimator precision (\Cref{CaseStudy}). This was achieved by
increasing the effective sample size from 2,179 to 4,505 valid steps (Table S1).
The improvement in estimator precision was weaker for F3-SH than for F3-S, as
the F3-SH model was more complex due to inclusion of additional interaction
terms. Point estimates for the habitat-selection and movement parameters were
similar for all 3 models, and evidence for scale dependency in habitat selection
was fairly weak. F3-S and F3-SH had similar AIC scores ($\Delta AIC \leq 1$;
Table S1), and the interaction terms were statistically significant only for the
step duration of 6 hours and only for one of the habitat covariates (Table S1).

\begin{figure}
  \begin{center}
  \includegraphics[width = \textwidth]{99_CaseStudy.png}
  \caption{Model results from the case study using GPS data collected on Apollo.
  In F1, forgiveness was set to one (only 2-hour steps were considered), whereas
  in F3-S and F3-SH a forgiveness of three was employed (allowing for step
  durations of up to 6 hours). In model F3-S, the step duration was interacted
  with step descriptors. In model F3-SH, step duration was interacted with step
  descriptors and habitat covariates. The bars indicate the 90\%, 95\%, and 99\%
  confidence intervals. Note that for simplicity, we omitted interactions with
  the step duration from this figure.}
  \label{CaseStudy}
  \end{center}
\end{figure}

\section{Discussion}
%Notably, both Munden et al.'s and McClintock's methods can
%accommodate step durations that are entirely random, as is often the case in
%marine applications, where locations can only be obtained when animals surface
%\citep{Costa.2010, Albertsen.2019}. Our focus, however, is on applications
%where data are schedule to be collected at regular intervals, but where
%observations are occasionally missing due to failures of the telemetry device
%to obtain a location. Nevertheless, Munden et al.'s approach motivates another
%solution we coin the \textit{na\"ive} approach.

We conducted a simulation study with known habitat and movement parameters to
investigate if retaining irregular animal locations via increased forgiveness
improves or worsens parameter estimation in iSSFs. We also tested the
performance of four different approaches that attempt to correct for potential
biases introduced by using temporally irregular data, and we compared them to an
uncorrected baseline approach. Our results demonstrated that retaining irregular
animal locations can improve the precision of estimators of habitat-selection
parameters but may lead to biased estimators of the parameters in
the movement kernel. Overall, our results highlight the potential benefits of
leveraging irregular animal locations, especially if an appropriate method for
handling irregular data is chosen.

The \textit{uncorrected} baseline approach ignored the fact that increasing
forgiveness introduced irregularity in the data. Consequently, estimators of the
parameters in the movement kernel were increasingly biased as forgiveness
increased due to the inclusion of steps with varying step durations when fitting
the model. Steps with longer step durations tended to have larger step lengths
and less directed turning angles, which led to an overestimation of $\theta$ and
underestimation of $k$ and $\kappa$. Yet, estimators of the habitat-selection
parameters remained unbiased or nearly so, even at high levels of forgiveness,
and they were more precise than the standard estimator (represented by the
\textit{uncorrected} approach with forgiveness = 1). These results highlight the
potential benefits that can be reaped by including additional data.

Similarly, the \textit{na\"ive} approach performed well when estimating
habitat-selection parameters but resulted in biased estimators of movement
parameters, especially for large forgiveness values. This result was
unsurprising given that our simulated trajectories were tortuous, and therefore,
step lengths were not linearly related to step durations. Indeed, we found that,
although there was a near linear relationship between step duration and the
(tentative) scale parameter ($\theta_0$), the relationships between step
duration and the (tentative) movement parameters $\kappa_0$ and $k_0$ were
non-linear (Figure S4). Overall, the usefulness of this approach appears highly
limited, as it is often not clear by what factor distributional parameters for
step lengths and turning angles should be multiplied to match the observed step
duration.

The \textit{dynamic+model} approach provided a flexible, easily implementable,
and powerful framework for retrieving precise and unbiased estimators of the
step-length and habitat-selection parameters, irrespective of the forgiveness
level. The estimator of the concentration parameter exhibited some bias, but
less than when using the \textit{uncorrected} and \textit{na\"ive} approaches.
To implement the \textit{dynamic+model} approach, we included interactions
between step descriptors (sl, log(sl), and cos(ta)) and step duration in the
conditional logistic regression model. This allowed the parameters of the
movement kernel to depend on the step duration. A complication, however, is that
turning angles are influenced by the step duration of both the current and
previous step (Figure S3). The bias in the concentration parameter likely arose
from only accounting for the step duration associated with the current step and
not the previous one. Moreover, fitting tentative distributions for different
step durations can be challenging due to some step durations occurring only
rarely. However, by resampling observed animal locations to different step
durations using the \texttt{track\_resample} function from the \texttt{amt}
\texttt{R}-package \citep{Signer.2019} the needed data can easily be generated.
We included step duration as a categorical covariate, yet there may be times
when it would be advantageous to treat it as a continuous covariate (e.g., with
its effect modeled using a low-degree polynomial or regression spline with few
degrees of freedom). Treating step duration as a continuous variable may help to
alleviate convergence issues in cases where some step durations are rare, and it
might allow applying the \textit{dynamic+model} approach to data that are
entirely irregular.

The \textit{multistep} approach also performed well and was relatively easy to
implement. This approach is somewhat \textit{ad hoc} in that it uses the
tentative movement parameters to generate random steps to match observed steps
with longer step durations (in multiples of $\Delta t$). It is similar to, but
slightly less principled, than the approach developed by \citet{Vales.2022},
which formally constructs the likelihood for multistep durations by integrating
out the missing steps. An advantage of this latter approach is that one can also
attempt to account for non-random missingness by explicitly modeling factors
related to the probability of obtaining a successful location
\citep{Vales.2022}. Nonetheless, integrating over the missing steps, as in
\citet{Vales.2022}, can be computationally intensive and prohibitive with large
data sets. Another downside of both of these approaches (the \textit{multistep}
approach and the approach of \citealp{Vales.2022}) is that they can only be
applied in cases where step durations are a fixed multiple of the regular
step duration; i.e., unlike the \textit{dynamic+model} approach, they cannot be
applied when data are highly irregular.

Of the methods we considered, the \textit{imputation} approach performed the
worst. It resulted in biased estimators of parameters in both the
habitat-selection function and the movement kernel. This bias likely resulted
from using an overly simplistic movement model to impute missing fixes.
Moreover, the imputation procedure may have led to imputed animal locations that
masked important selection properties, therefore leading to inaccurate parameter
estimates. While this approach appears to perform well with hidden Markov
movement models \citep{McClintock.2017}, we advise against its use with iSSFs.

For the scenarios we considered in our simulation study, the estimators of
habitat-selection parameters were insensitive to the inclusion of irregular data
and performed well, except for the \textit{imputation} approach. This suggests
that accounting for irregular step durations may not be particularly important
if one is only interested in the habitat-selection function. When the movement
kernel is also of interest, we suggest the \textit{dynamic+model} approach,
since it is flexible, easy to implement, and allows one to use more data than
the traditional approach that requires bursts of regular data, leading to more
precise estimators.

Several authors have emphasized that movement and habitat-selection parameters
in an SSF are scale dependent and should be expected to change as the sampling
frequency changes \citep[see for example][]{Avgar.2016, Signer.2017,
Fieberg.2021}. Furthermore, \citet{Barnett.2008} developed an analytical
framework for investigating scale dependence and showed that habitat-selection
parameters should depend on the relative width of the movement kernel in
relation to habitat heterogeneity. Thus, the relative insensitivity of the
habitat-selection parameters to the inclusion of steps with varying step
duration was somewhat unexpected. It would be interesting to explore the
robustness of this result across a wider range of simulation scenarios in the
future.

More generally, the spatial scale of a habitat-selection analysis has been
recognized as an important factor, which is why \citet{Johnson.1980} proposed a
hierarchical framework for examining habitat-selection across different orders
(e.g., species range, individual home range, within a home range). Johnson's
proposed framework acknowledges that habitat-selection may act differently at
different scales, and that the interpretation of ecological processes changes
depending on the spatial scale at which they are investigated \citep{Wiens.1989,
Levin.1992}. This understanding has encouraged scientists to conduct extensive
scaling analyses and to comprehensively examine habitat-selection at multiple
scales \citep{DeCesare.2012, McGarigal.2016, Pitman.2017, Zeller.2017}. In
studies employing SSAs, the issue of scale is often neglected, and data are most
frequently analyzed at the spatio-temporal scale at which they were collected.
This choice maximizes the number of locations that can be used in the analysis,
yet prevents a thorough understanding of scale dependency. The use of irregular
data in SSAs poses another challenge, as steps with unequal step durations may
reflect selection processes occurring at different scales. The severity of this
issue obviously depends on the original sampling frequency, the degree of
missingness, and the scale at which animals are making decisions that are
relevant in terms of their movement behavior and habitat selection. By including
irregular animal locations via increased forgiveness, we may therefore average
over selection processes occurring at multiple scales, which could produce
estimates of habitat-selection parameters that are misleading due to
contradictory effect signs at different scales. To better account for such
scale-dependent processes, it may be beneficial to include interactions between
step duration and habitat features (e.g., \textsf{dist}, \textsf{elev},
\textsf{forest}), thus allowing habitat-selection parameters to also vary as a
function of step duration. We demonstrated how this could be implemented in the
case study.

It is important to note that we considered a limited number of scenarios in our
simulation study. For instance, we assumed that animal locations were missing at
random, i.e., failure to obtain a fix was unrelated to habitat types, time of
the day, etc. However, several studies have shown that missingness is often
non-random and related to difficulties with satellite transmission due to
topography \citep{Lewis.2007}, canopy cover \citep{Phillips.1998, DeCesare.2005,
Hansen.2008}, time of the day \citep{Graves.2006}, animal behavior
\citep{Mattisson.2010}, or collar orientation \citep{Deon.2005}. In fact,
\citet{Vales.2022} highlighted that missingness and the associated
under-representation of certain habitat types may lead to biased estimators of
parameters in iSSFs, but that accounting for the probability of obtaining a
location in differing environmental conditions may alleviate this bias. Future
studies should strive to further investigate these relationships and examine how
our proposed approaches perform when missingness is habitat-dependent.

A major benefit of using iSSFs is the ability to allow an individual's movement
kernel to depend on local habitat features \citep{Avgar.2016}. In our simulation
study, we considered simplified scenarios where the movement kernel was
unchanging, which simplified the simulation and inference. Nevertheless, such
interactions often play a crucial role in real ecosystems. For instance,
\citet{Dickie.2020} employed iSSFs and revealed that several large mammal
species moved faster while on linear features. Similarly, \citet{Hofmann.2023}
found that African wild dogs moved significantly slower and less directed in
areas that were covered by floodwater. Future studies could investigate
simulation scenarios in which individuals alter their movement tendencies in
response to local environmental features (i.e., models with habitat dependent
movement kernels) and examine how this influences the robustness of our proposed
approaches.

While our results suggest that irregularity due to missing animal locations can
effectively be accounted for in iSSFs and that increasing the forgiveness, thus
allowing for inclusion of irregular data, improves estimator precision, we also
found a decreasing marginal benefit of increased forgiveness. In fact,
increasing the forgiveness beyond a value of two (i.e., allowing for steps of
twice the regular step duration) only marginally improved model performance in
our case. This can also be seen in \Cref{NumberOfSteps}, which shows that the
largest number of steps that can be gained is when increasing the forgiveness
level from one to two. Having a higher forgiveness beyond two may thus not even
be necessary, therefore limiting the need to correct biases emerging from the
inclusion of irregular data.

Although we focused on the case of missing location data, the proposed
approaches may also prove useful for situations where sampling is irregular for
other reasons. For example, it is not uncommon to adjust sampling regimes after
a preliminary phase, following improvements to collar-battery-lifetime, or for
sampling rates to vary depending on type and manufacturer of the collar
\citep{Brown.2023}. Similarly, it is common practice to adjust the GPS regime to
the biology of the focal species and only record data during a specific time of
the day (e.g., \citealp{Broekhuis.2013, Cozzi.2013, Elliot.2014}). These
irregularities might be addressed using the \textit{dynamic+model} approach,
with interactions between step duration and movement descriptors. Interactions
between step duration and habitat-selection covariates should also be
considered, particularly if the sampling regime is adjusted to coincide with
changes in animal behavior. We expect this approach will work fairly well in
many cases, but we might expect a slight bias in the estimated concentration
parameters, as observed in our simulation study.

Our study contributes to the growing body of literature that extends iSSFs and
improves the method's robustness under various conditions. This includes
approaches for modeling irregular data \citep{Munden.2021, Eisaguirre.2024},
accounting for spatial dependence among residuals \citep{ArceGuillen.2023},
methodological frameworks for fitting iSSFs with random slopes \citep{Muff.2020}
and random smooths \citep{Klappstein.2024}, incorporating the probability of
successfully obtaining an animal location in different habitat conditions
\citep{Vales.2022}, and considering the behavioral states of the tracked animals
\citep{Klappstein.2023, Pohle.2023}.

In conclusion, our study shows that inclusion of irregular animal locations can
improve model performance, yet only when an appropriate approach to account for
irregularity is selected. Here, the \textit{dynamic+model} and
\textit{multistep} approaches performed well and resulted in improved estimators
of habitat-selection and movement parameters, even at elevated levels of
missingness and forgiveness. Both methods are easy to implement, and the
associated models can readily be fitted using the R-packages \texttt{amt}
\citep{Signer.2019}, \texttt{survival} \citep{Therneau.2023}, \texttt{coxme}
\citep{Therneau.2022}, and \texttt{mgcv} \citep{Wood.2011, Wood.2017,
Klappstein.2024}. To facilitate uptake and encourage use of the proposed
approaches among practitioners, we provide all of our codes through an online
repository, which includes an example application of the \textit{dynamic+model}
approach. With this, we hope practitioners will rethink the common use of
discarding large portions of data and instead use methods that can accommodate
irregular data.

\section{Authors' Contributions}

D.D.H., G.C., and J.F. conceived the study and designed methodology; D.D.H.
implemented the analysis, J.F. assisted with modeling; D.D.H. wrote the first
draft of the manuscript and all authors contributed to the drafts at several
stages and gave final approval for publication.

\section{Data Availability}

Code to reproduce this study is available through the University of Minnesota's
Data repository (\url{https://doi.org/10.13020/6wcd-6s43},
\citealp{Hofmann.2023a}). The repository also includes two example analyses that
showcase the application of the \textit{dynamic+model} approach to simulated
data, as well as to the GPS data of Apollo.

\section{Acknowledgements \& Funding}

D.D.H. was funded by a Swiss National Science Foundation Grant No.
310030\_204478 rewarded to G.C. J.F. was supported by the National Aeronautics
and Space Administration award 80NSSC21K1182 and received partial salary support
from the Minnesota Agricultural Experimental Station. We thank two anonymous
reviewers for their constructive and detailed feedback, which helped to improve
the quality of this manuscript.

\newpage
\begingroup
\singlespacing
\ifthenelse{\boolean{usebiblatex}}{
  \begin{refcontext}[sorting=nyt]
  \printbibliography
  \end{refcontext}
}{
  \bibliography{LiteratureBibtex}
}
\endgroup
\end{document}
